% !TeX root = ../main.tex
\section{Introduction}
\label{sec:introduction}

Compute-in-memory (CIM) architectures execute dense matrix-vector multiplications directly within memory arrays, eliminating the data-movement bottleneck that limits conventional digital systems \citep{horowitz2014computing,peng2020dnnneurosim,lanza2025memristor}. Across emerging analog memories, including resistive memories, phase-change memory (PCM), and organic optoelectronic devices, deployment accuracy is shaped by finite conductance precision, device-to-device variability, cycle-to-cycle noise, and retention drift \citep{rasch2021aihwkit,crosssim2024,xu2025emerging,guo2024organic,photonics2025organicreview}. These effects are often characterized at the device level, but their system-level consequences for modern vision backbones remain difficult to rank before fabrication.

However, translating device-level models into deployment-facing model accuracy faces a critical methodological gap. Mature CIM simulators expose increasingly detailed analog effects \citep{peng2020dnnneurosim,lammie2022memtorch,rasch2021aihwkit,crosssim2024}, yet the algorithmic and physical failure modes are often entangled: a low-precision model can fail because hardware-aware training overfits one fixed mismatch realization, because quantization is too coarse, or because a realistic memory substrate drifts after programming. Device-specific studies, including organic and optoelectronic demonstrations, further highlight this gap because variability, ADC resolution, retention-induced drift, and transfer across fabricated instances are frequently absent from network-level analyses or treated only qualitatively \citep{zeng2023organicmemristor,jung2024organicfilaments,lin2016physical}. A deployment workflow therefore needs to separate algorithmic robustness from physical precision-retention limits.

Here we present a behavioral simulation framework that evaluates analog neural network deployment under both idealized pure-quantization and realistic device-physics regimes. We apply a mixed-signal deployment model to the Tiny-ViT-5M backbone to systematically disentangle algorithmic generalization failures from physical retention limits. Three concrete contributions follow. First, we expose a fresh-instance transfer failure in standard fixed-instance hardware-aware training at extreme low precision (4-bit), where accuracy collapses to near-chance. We introduce Ensemble Hardware-Aware Training (HAT), which resamples device-to-device mismatch masks each epoch, rescuing the 4-bit pure-quantization regime and restoring fresh-instance accuracy to 86.16$\pm$0.19\%. Second, we translate this algorithmic rescue into a realistic phase-change memory (PCM) simulation regime and observe stable convergence under the tested PCM UnitCell setting where pure quantization fails. Third, we map a 4/6/8-bit precision-retention deployment frontier for PCM. In this tested sweep, 4-bit is trainable but drift-limited, 8-bit is drift-flat, and 6-bit provides the best observed Pareto midpoint between model compression and long-term inference reliability.

The framework is intentionally behavioral rather than circuit-accurate, focusing on macroscopic precision and drift phenomena rather than array-level parasitics (IR drop, sneak paths) which remain crucial future considerations. Additional supporting analyses, including front-end compensation and zero-shot transfer to alternative organic optoelectronic profiles, are provided in the Supplementary Information. Its purpose is to provide a reproducible workflow for comparing mitigation strategies and establishing physical deployment frontiers.
